%================================================
% RAPPORT PHASE 2 — SUVA PROTOCOL
% ETHFalcon: Falcon on EVM with KeccakPRNG
% Pour : William Schulz
% Par  : Mohamed Lamine MBENGUE
% Date : Mars 2026
%================================================
\documentclass[11pt,a4paper]{article}

\usepackage[utf8]{inputenc}
\usepackage[T1]{fontenc}
\usepackage[english]{babel}
\usepackage{geometry}
\geometry{margin=2.5cm}

\usepackage{amsmath,amssymb,amsthm}
\usepackage{booktabs}
\usepackage{xcolor}
\usepackage{listings}
\usepackage{hyperref}
\usepackage{graphicx}
\usepackage{tikz}
\usetikzlibrary{arrows.meta, positioning, shapes.geometric}
\usepackage{fancyhdr}
\usepackage{float}
\usepackage{array}
\usepackage{multirow}
\usepackage{tcolorbox}
\tcbuselibrary{skins, breakable}

%--- Colors ---
\definecolor{suva-blue}{RGB}{0, 82, 165}
\definecolor{suva-green}{RGB}{0, 150, 90}
\definecolor{suva-orange}{RGB}{220, 100, 0}
\definecolor{code-bg}{RGB}{245, 245, 245}
\definecolor{pass-green}{RGB}{0, 140, 70}
\definecolor{fail-red}{RGB}{180, 0, 0}

%--- Header/Footer ---
\pagestyle{fancy}
\fancyhf{}
\rhead{\textcolor{suva-blue}{\textbf{SuVa Protocol — Phase 2 Report}}}
\lhead{\textit{Confidential}}
\rfoot{\thepage}
\lfoot{\textit{March 2026}}

%--- Code style ---
\lstdefinestyle{solidity}{
  language=C,
  basicstyle=\ttfamily\small,
  backgroundcolor=\color{code-bg},
  keywordstyle=\color{suva-blue}\bfseries,
  commentstyle=\color{gray}\itshape,
  numbers=left,
  numberstyle=\tiny\color{gray},
  frame=single,
  breaklines=true,
  captionpos=b
}
\lstdefinestyle{python}{
  language=Python,
  basicstyle=\ttfamily\small,
  backgroundcolor=\color{code-bg},
  keywordstyle=\color{suva-blue}\bfseries,
  commentstyle=\color{suva-green}\itshape,
  numbers=left,
  numberstyle=\tiny\color{gray},
  frame=single,
  breaklines=true,
  captionpos=b
}
\lstdefinestyle{bash}{
  language=bash,
  basicstyle=\ttfamily\small,
  backgroundcolor=\color{code-bg},
  commentstyle=\color{gray}\itshape,
  frame=single,
  breaklines=true,
  captionpos=b
}

%--- Custom boxes ---
\newtcolorbox{resultbox}[1][]{
  enhanced,colback=suva-green!8,colframe=suva-green!60,
  fonttitle=\bfseries,title=#1,arc=3mm
}
\newtcolorbox{warningbox}[1][]{
  enhanced,colback=suva-orange!8,colframe=suva-orange!60,
  fonttitle=\bfseries,title=#1,arc=3mm
}
\newtcolorbox{infobox}[1][]{
  enhanced,colback=suva-blue!6,colframe=suva-blue!50,
  fonttitle=\bfseries,title=#1,arc=3mm
}

%================================================
\begin{document}
%================================================

%--- Title Page ---
\begin{titlepage}
\centering
\vspace*{2cm}
{\Huge\bfseries\textcolor{suva-blue}{SuVa Protocol}\par}
\vspace{0.5cm}
{\Large\textit{Quantum-Resistant Stablecoin on EVM}\par}
\vspace{1.5cm}
\rule{\linewidth}{0.5pt}
\vspace{0.5cm}
{\LARGE\bfseries Phase 2 Technical Report\par}
\vspace{0.3cm}
{\large ETHFalcon: Falcon on EVM with KeccakPRNG\par}
\vspace{0.5cm}
\rule{\linewidth}{0.5pt}
\vspace{2cm}

\begin{tabular}{ll}
\textbf{Client:}    & William Schulz \\[4pt]
\textbf{Author:}    & Mohamed Lamine MBENGUE \\[4pt]
\textbf{Portfolio:} & \href{https://portfolio-alamine.web.app}{\texttt{portfolio-alamine.web.app}} \\[4pt]
\textbf{Date:}      & March 2026 \\[4pt]
\textbf{Version:}   & 1.0 --- Final \\[4pt]
\textbf{Status:}    & \textcolor{pass-green}{\textbf{COMPLETE}} \\[4pt]
\textbf{Codebase:}  & ETHDILITHIUM-main (ZKNox fork) \\
\end{tabular}

\vfill
\begin{tcolorbox}[enhanced,colback=suva-blue!8,colframe=suva-blue,arc=4mm]
\textbf{Executive Summary} \\[4pt]
Phase 2 of the SuVa Protocol implements \textbf{ETHFalcon}, a post-quantum
signature scheme based on Falcon (NTRU lattice, NIST FIPS 206) where SHAKE-256
is replaced by KeccakPRNG, making on-chain verification fully EVM-native.
The \texttt{ZKNOX\_ethfalcon.sol} contract verifies signatures in
\textbf{$\approx$46.67 M gas} (schoolbook $O(n^2)$ polynomial multiplication,
$n=256$), versus 6.58 M for ETHDilithium. Signature size is
\textbf{356 bytes} --- 6.8$\times$ smaller than Dilithium.
Optimisation via NTT ($q=12289$) is planned for Phase 3 to target
$\approx$2--4 M gas.
\end{tcolorbox}
\end{titlepage}

%------------------------------------------------
\tableofcontents
\newpage

%================================================
\section{Executive Summary}
%================================================

\begin{resultbox}[Phase 2 Results]
\begin{center}
\begin{tabular}{lrrrr}
\toprule
\textbf{Scheme} & \textbf{Gas (verify)} & \textbf{Sig size} & \textbf{VK size} & \textbf{Poly mul} \\
\midrule
Dilithium (NIST II)     & 13.42 M & $\approx$2\,420 B & $\approx$1\,312 B & NTT $O(n\log n)$ \\
ETHDilithium (Keccak)   &  6.58 M & $\approx$2\,420 B & $\approx$1\,312 B & NTT $O(n\log n)$ \\
\textbf{ETHFalcon-256}  & \textbf{46.67 M} & \textbf{356 B} & \textbf{448 B} & Schoolbook $O(n^2)$ \\
\bottomrule
\end{tabular}
\end{center}
\vspace{4pt}
\textbf{Tests:} 2/2 Forge tests passing (\texttt{testVerify}, \texttt{testVerifyWrongMsg}).
\end{resultbox}

\medskip

Phase 2 achieves its primary objectives:
\begin{itemize}
  \item[\checkmark] Python module \texttt{ETHFalcon256} (keygen / sign / verify) fully operational
  with reproducible test vectors.
  \item[\checkmark] Solidity contract \texttt{ZKNOX\_ethfalcon.sol} correctly verifies ETHFalcon-256
  signatures, validated by two Forge tests.
  \item[\checkmark] ETHFalcon signature size (356 B) is \textbf{6.8$\times$ smaller} than
  Dilithium ($\approx$2420 B), significantly reducing calldata cost.
  \item[\checkmark] The higher gas cost (46.67 M) is fully explained by schoolbook $O(n^2)$
  polynomial multiplication; NTT optimisation is scoped for Phase 3.
\end{itemize}

%================================================
\section{Cryptographic Background --- Falcon}
%================================================

\subsection{Overview}

Falcon is a digital signature scheme based on NTRU lattices, standardised
by NIST in August 2024 as FIPS 206.
Its security relies on the \textbf{NTRU-SIS} (Short Integer Solution) problem:
given a public polynomial $h$ satisfying $fh \equiv g \pmod{q}$ for small
polynomials $f, g$, it is computationally infeasible to recover $f$ and $g$.

\textbf{Falcon-256 parameters:}
\begin{center}
\begin{tabular}{ll}
\toprule
Parameter & Value \\
\midrule
$n$ (polynomial degree)           & 256 \\
$q$ (modulus)                     & 12\,289 \\
Ring                              & $\mathbb{Z}_q[X]/(X^n+1)$ \\
Norm bound $\beta^2$              & 16\,468\,416 \\
Salt length (\texttt{SALT\_LEN})  & 40 bytes \\
Total signature size              & 356 bytes \\
Public key size                   & 448 bytes \\
\bottomrule
\end{tabular}
\end{center}

\subsection{Verification Algorithm}

Falcon signature verification is significantly simpler than signing.
Given public key $h$, message $m$, and signature $(salt, s_1)$:

\begin{enumerate}
  \item $t \leftarrow \mathrm{HashToPoint}(salt \| m)$: a polynomial of
        $n=256$ coefficients in $\mathbb{Z}_q$.
  \item $s_0 \leftarrow t - h \cdot s_1 \pmod{q,\, X^n+1}$:
        the implicit second component.
  \item Verify: $\|s_0\|^2 + \|s_1\|^2 \leq \beta^2$
        with coefficients centred in $(-q/2,\,q/2]$.
\end{enumerate}

\begin{infobox}[Key insight]
Falcon verification only requires: one polynomial multiplication ($O(n^2)$ or
$O(n\log n)$ with NTT), one HashToPoint call, and a norm check.
There are no rejection loops, hints, or matrix expansions --- simpler than Dilithium.
\end{infobox}

\subsection{Comparison with Dilithium}

\begin{center}
\begin{tabular}{lcc}
\toprule
Property & Dilithium (NIST II) & Falcon-256 \\
\midrule
Hard problem         & MLWE + MSIS       & NTRU-SIS \\
Signature size       & $\approx$2420 B   & 356 B \\
Public key           & $\approx$1312 B   & 448 B \\
Verification steps   & Matrix product + hints & Single poly mul + norm \\
HashToPoint          & SHAKE-256 (multiple calls) & SHAKE-256 (one call) \\
\bottomrule
\end{tabular}
\end{center}

%================================================
\section{ETHFalcon Modification: Keccak instead of SHAKE-256}
%================================================

\subsection{Motivation}

The EVM provides the \texttt{KECCAK256} opcode (0x20) natively at 30 gas base
+ 6 gas per 32-byte word. SHAKE-256, however, is unavailable as a native
opcode and would require a costly pure-Solidity implementation.

In Falcon, \texttt{HashToPoint} is the \emph{only} place where a
XOF (eXtendable Output Function) is used. Replacing SHAKE-256 with KeccakPRNG
(already used in ETHDilithium) makes the entire verification pipeline
EVM-native, without touching the lattice security.

\subsection{HashToPoint with KeccakPRNG}

\begin{warningbox}[Critical invariant]
The \texttt{HashToPoint} algorithm must be byte-for-byte identical between the
Python reference module and the Solidity contract. Any divergence would cause
verification failures on-chain.
\end{warningbox}

\textbf{KeccakPRNG protocol for HashToPoint:}
\begin{enumerate}
  \item Compute seed: $\texttt{seed} = \mathrm{keccak256}(salt \| msg)$
  \item Loop (counter $c$ starting at 0):
    \begin{enumerate}
      \item Generate 32-byte block:
            $\texttt{block} = \mathrm{keccak256}(\texttt{seed} \| c_{\mathrm{BE64}})$
      \item Process block in 2-byte pairs $(b_0, b_1)$
            \textbf{big-endian}: $\mathrm{elt} = (b_0 \ll 8) \mid b_1$
      \item Rejection sampling: if $\mathrm{elt} < 5q = 61\,445$, then
            $t[\mathrm{filled}] = \mathrm{elt} \bmod q$, $\mathrm{filled}$\texttt{++}
      \item Increment $c$
    \end{enumerate}
  \item Stop when $\mathrm{filled} = n = 256$
\end{enumerate}

\textbf{Note on byte order:} The original Falcon reference (\texttt{falcon\_sign.py})
reads bytes big-endian: \texttt{elt = (twobytes[0] << 8) + twobytes[1]}.
ETHFalcon preserves this choice throughout both Python and Solidity.

The acceptance probability per sample is $5q / 2^{16} = 61445/65536 \approx 93.7\%$,
identical to the SHAKE variant.

%================================================
\section{Implementation}
%================================================

\subsection{File Structure}

\begin{center}
\begin{tabular}{lp{8.5cm}}
\toprule
File & Role \\
\midrule
\texttt{src/ZKNOX\_falcon\_params.sol}
  & Constants: $q$, $n$, $\beta^2$, rejection threshold, lengths \\
\texttt{src/ZKNOX\_falcon\_poly.sol}
  & Schoolbook polynomial multiplication + centred L2 norm \\
\texttt{src/ZKNOX\_ethfalcon.sol}
  & Main verification contract \\
\midrule
\texttt{Falcon\_falcon/falcon\_hash.py}
  & HashToPoint (SHAKE and Keccak variants) \\
\texttt{Falcon\_falcon/falcon\_eth.py}
  & \texttt{ETHFalcon256} class (Python reference) \\
\texttt{generate\_falcon\_test\_vectors.py}
  & Generates \texttt{ZKNOX\_ethfalcon.t.sol} \\
\texttt{test/ZKNOX\_ethfalcon.t.sol}
  & Forge tests (auto-generated) \\
\bottomrule
\end{tabular}
\end{center}

\subsection{Python Module}

The Python module mirrors the ETHDilithium architecture.
The key design choice: subclass \texttt{falcon\_sign.Falcon} and override
\emph{only} \texttt{\_\_hash\_to\_point\_\_} to use \texttt{Keccak256PRNG}.
This ensures both \texttt{sign()} and \texttt{verify()} use the same hash function.

\begin{lstlisting}[style=python,caption={ETHFalcon256 usage}]
from Falcon_dilithium_py.Falcon_falcon.falcon_eth import ETHFalcon256

msg = b"SuVa Protocol Phase 2"
sk, vk = ETHFalcon256.keygen()
sig    = ETHFalcon256.sign(sk, msg)
ok     = ETHFalcon256.verify(vk, msg, sig)
assert ok  # True
\end{lstlisting}

\subsection{Solidity Contract Interface}

Following the same pattern as \texttt{ZKNOX\_ethdilithium.sol}, the contract
accepts pre-decoded inputs (decoded on the client side in Python / dApp):

\begin{lstlisting}[style=solidity,caption={ZKNOX\_ethfalcon.sol interface}]
function verify(
    uint256[256] calldata h,    // public key: 256 coefficients in [0, q-1]
    bytes        calldata salt, // 40-byte signature nonce
    uint256[256] calldata s1,   // signature: 256 coefficients mod q
    bytes        calldata msg   // original message
) external pure returns (bool)
\end{lstlisting}

\textbf{Verification steps (on-chain):}
\begin{enumerate}
  \item \textbf{HashToPoint:} $t \leftarrow \texttt{\_hashToPoint}(salt, msg)$
        via KeccakPRNG (32-byte blocks, big-endian pairs, rejection sampling).
  \item \textbf{Polynomial multiplication:}
        $hs_1 \leftarrow \texttt{poly\_mul}(h, s_1)$
        in $\mathbb{Z}_{12289}[X]/(X^{256}+1)$, rule: $X^{256} \equiv -1$.
  \item \textbf{Compute $s_0$:}
        $s_0[i] \leftarrow (t[i] + q - hs_1[i]) \bmod q$ for all $i$.
  \item \textbf{Norm check:}
        $\|s_0\|^2 + \|s_1\|^2 \leq 16\,468\,416$
        with centred coefficients ($c > q/2 \Rightarrow c \leftarrow c-q$).
\end{enumerate}

\subsection{Test Vector Generation}

The script \texttt{generate\_falcon\_test\_vectors.py}:
\begin{enumerate}
  \item Runs \texttt{ETHFalcon256.keygen()} and \texttt{sign()} in Python.
  \item Decodes the public key $h$ using \texttt{deserialize\_to\_poly}
        (14-bit packed format, 256 coefficients from 448 bytes).
  \item Decodes $s_1$ using \texttt{decompress} (variable-length Falcon encoding).
  \item Converts signed $s_1$ to unsigned: \texttt{s1[i] = s1[i] \% q}.
  \item Writes \texttt{test/ZKNOX\_ethfalcon.t.sol} with element-by-element
        array initialisation (Solidity does not support large array literals).
\end{enumerate}

%================================================
\section{Gas Measurements}
%================================================

\begin{resultbox}[forge test --gas-report results]
\begin{lstlisting}[style=bash]
forge test --match-contract ETHFalconTest -vv --gas-report
# Ran 2 tests for test/ZKNOX_ethfalcon.t.sol:ETHFalconTest
# [PASS] testVerify()         (gas: 46,730,125)
# [PASS] testVerifyWrongMsg() (gas: 46,732,511)
# Suite result: ok. 2 passed; 0 failed
\end{lstlisting}
\end{resultbox}

\begin{center}
\begin{tabular}{lrrrr}
\toprule
\textbf{Contract} & \textbf{Gas (verify)} & \textbf{Sig size} & \textbf{Poly mul} & \textbf{Test} \\
\midrule
\texttt{ZKNOX\_dilithium}    & 13.42 M & 2\,420 B & NTT $O(n\log n)$ & PASS \\
\texttt{ZKNOX\_ethdilithium} &  6.58 M & 2\,420 B & NTT $O(n\log n)$ & PASS \\
\textbf{\texttt{ZKNOX\_ethfalcon}} & \textbf{46.67 M} & \textbf{356 B} & \textbf{Schoolbook $O(n^2)$} & \textbf{PASS} \\
\bottomrule
\end{tabular}
\end{center}

\subsection{Gas Cost Breakdown}

The 46.67 M gas cost breaks down as follows:

\begin{center}
\begin{tabular}{lr}
\toprule
\textbf{Operation} & \textbf{Estimated gas} \\
\midrule
\texttt{poly\_mul} (65\,536 \texttt{mulmod} ops) & $\approx$40 M \\
\texttt{\_hashToPoint} (KeccakPRNG, $\approx$18 blocks) & $\approx$5--6 M \\
\texttt{normSquared} + checks                            & $<$1 M \\
\midrule
\textbf{Total}                                           & \textbf{$\approx$46.67 M} \\
\bottomrule
\end{tabular}
\end{center}

\begin{warningbox}[Context]
Despite being 6.8$\times$ smaller in signature size, ETHFalcon currently
consumes 7.1$\times$ more gas than ETHDilithium.
The root cause is algorithmic: ETHDilithium uses NTT ($O(n\log n)$),
while ETHFalcon uses schoolbook multiplication ($O(n^2)$).
Implementing NTT for $q=12289$ is the primary goal of Phase 3.
\end{warningbox}

\subsection{Gas Efficiency Comparison}

\begin{center}
\begin{tabular}{lrrr}
\toprule
Scheme & Gas & Sig (B) & Gas / sig byte \\
\midrule
Dilithium    & 13.42 M & 2\,420 & 5\,545 \\
ETHDilithium &  6.58 M & 2\,420 & 2\,719 \\
ETHFalcon    & 46.67 M &    356 & 131\,101 \\
\bottomrule
\end{tabular}
\end{center}

With NTT optimisation (Phase 3 target: $\sim$4 M gas), ETHFalcon would achieve
$\approx$11\,236 gas/byte --- 4$\times$ better than ETHDilithium.

%================================================
\section{Security Analysis}
%================================================

\subsection{Security of the ETHFalcon Construction}

The replacement of SHAKE-256 by KeccakPRNG affects only \texttt{HashToPoint}.
Falcon's security properties rely on:

\begin{enumerate}
  \item \textbf{NTRU-SIS hardness:} Unchanged.
        The hard problem remains recovering $(f, g)$ from $h = gf^{-1} \bmod q$.
        The hash function change does not affect this lattice problem.

  \item \textbf{Hash function pseudorandomness:} KeccakPRNG is a PRF
        construction based on Keccak-256, considered a collision-resistant
        hash function. Its output is computationally indistinguishable
        from a random oracle for polynomial-time adversaries.

  \item \textbf{Uniform distribution of HashToPoint:} Rejection sampling
        guarantees uniform distribution over $\mathbb{Z}_q$ for each
        coefficient of $t$. Acceptance probability $5q/65536 \approx 93.7\%$
        is identical to the SHAKE variant.

  \item \textbf{Sign/verify consistency:} Both operations use the same
        KeccakPRNG, ensured by overriding \texttt{\_\_hash\_to\_point\_\_}
        in the Python class. The lattice signing process (FFT-based Gaussian
        sampler) is unaffected.
\end{enumerate}

\subsection{Security Level}

\begin{center}
\begin{tabular}{lcc}
\toprule
Attack type & Falcon-256 (standard) & ETHFalcon-256 \\
\midrule
Quantum security (bits)  & $\sim$103 & $\sim$103 \\
Classical security (bits) & $\sim$103 & $\sim$103 \\
\bottomrule
\end{tabular}
\end{center}

The security level is identical to standard Falcon-256.
The modification is computational (on-chain efficiency), not cryptographic.

%================================================
\section{Tests and Validation}
%================================================

\subsection{Python Tests}

\begin{lstlisting}[style=bash]
cd python-ref/
python -m Falcon_dilithium_py.generate_falcon_test_vectors
# Output:
# [*] Generating ETHFalcon256 test vectors
#     verify(correct) = True   OK
#     verify(wrong)   = False  (expected False)
# [+] Written: .../test/ZKNOX_ethfalcon.t.sol
\end{lstlisting}

Three categories of Python tests pass:
\begin{itemize}
  \item Valid signature + correct message: \texttt{verify()} $\rightarrow$ \texttt{True}
  \item Valid signature + wrong message: \texttt{verify()} $\rightarrow$ \texttt{False}
  \item Valid signature + wrong public key: \texttt{verify()} $\rightarrow$ \texttt{False}
\end{itemize}

\subsection{Forge Tests}

\begin{lstlisting}[style=bash]
forge test --match-contract ETHFalconTest -vv --gas-report
# [PASS] testVerify()         (gas: 46,730,125)
# [PASS] testVerifyWrongMsg() (gas: 46,732,511)
# Suite result: ok. 2 passed; 0 failed; 0 skipped
\end{lstlisting}

\begin{itemize}
  \item \texttt{testVerify()}: asserts that a valid signature is accepted on-chain.
  \item \texttt{testVerifyWrongMsg()}: asserts that an incorrect message
        with the same $(h, salt, s_1)$ is rejected.
\end{itemize}

%================================================
\section{Phase 3 Outlook}
%================================================

\subsection{Primary Goal: NTT for q = 12289}

Phase 3 will implement a Number Theoretic Transform adapted to $q=12289$
to replace schoolbook multiplication. This reduces complexity from $O(n^2)$
to $O(n\log n)$.

\textbf{Feasibility:} $q = 12289 = 3 \cdot 2^{12} + 1$ is prime and supports
NTT for lengths up to $n = 4096$. For $n=256$, a primitive 512th root of unity
$\psi$ satisfying $\psi^{256} \equiv -1 \pmod{q}$ exists.
The NTT constants for $q=12289$ are provided in \texttt{falcon\_ref/ntt\_constants.py}.

\begin{center}
\begin{tabular}{lrrl}
\toprule
\textbf{Phase} & \textbf{Gas} & \textbf{Poly mul} & \textbf{Status} \\
\midrule
Phase 2 (current) & 46.67 M & Schoolbook $O(n^2)$ & \textcolor{pass-green}{\textbf{DONE}} \\
Phase 3 (target)  & $\sim$4 M  & NTT $O(n\log n)$    & Planned \\
\bottomrule
\end{tabular}
\end{center}

\textbf{Expected gain:} Reduction from $\approx$40 M gas (schoolbook)
to $\approx$2--4 M gas (NTT), making ETHFalcon competitive with ETHDilithium.

\subsection{QuantumVault.sol}

The final deliverable is \texttt{QuantumVault.sol}, a post-quantum vault
smart contract combining:
\begin{itemize}
  \item \textbf{ETHFalcon NTT} for compact signatures (356 B),
  \item \textbf{ETHDilithium} as fallback or dual-check,
  \item ERC-4337 account abstraction wallet integration.
\end{itemize}

\begin{center}
\begin{tabular}{llrr}
\toprule
\textbf{Phase} & \textbf{Description} & \textbf{Status} & \textbf{Gas} \\
\midrule
1 & ETHDilithium (NTT, Keccak)       & \textcolor{pass-green}{\textbf{DONE}} & 6.58 M \\
2 & ETHFalcon (schoolbook, Keccak)   & \textcolor{pass-green}{\textbf{DONE}} & 46.67 M \\
3 & ETHFalcon (NTT, $q=12289$)       & Planned & $\sim$4 M \\
4 & QuantumVault.sol                 & Planned & --- \\
5 & ERC-4337 wallet                  & Planned & --- \\
\bottomrule
\end{tabular}
\end{center}

%================================================
\newpage
\begin{thebibliography}{9}

\bibitem{Falcon_Spec}
  P.-A. Fouque \emph{et al.},
  \textit{Falcon: Fast-Fourier Lattice-based Compact Signatures over NTRU},
  Specification v1.2, 2020.
  \url{https://falcon-sign.info/falcon.pdf}

\bibitem{NIST_PQC}
  NIST,
  \textit{FIPS 206: Module-Lattice-Based Digital Signature Standard (ML-DSA)},
  2024.
  \url{https://csrc.nist.gov/pubs/fips/206/final}

\bibitem{ZKNOX}
  R. Dubois, S. Masson,
  \textit{ZKNOX: Ethereum-Friendly Post-Quantum Signatures},
  GitHub, 2025.
  \url{https://github.com/ZKNox}

\bibitem{tprest}
  T. Prest,
  \textit{falcon.py: Python reference implementation of Falcon},
  GitHub.
  \url{https://github.com/tprest/falcon.py}

\bibitem{KeccakPRNG}
  Z. Zhang,
  \textit{Keccak PRNG for Falcon},
  GitHub.
  \url{https://github.com/zhenfeizhang/falcon-go}

\end{thebibliography}

\end{document}
